\section*{Bab \ref{ch:g9:linfunc} --- Fungsi Linear dan Fungsi Afin}

\begin{solution}{exo:g9:linfunc:1}
Linear: $f$ dan $k$ (berbentuk $ax$). Afin tetapi bukan linear: $g$
(\omterm{prop:g9:linfunc:affinegraph}{gradiennya} $3$, titik potongnya $-5$) dan $m$ (tetap, \omterm{prop:g9:linfunc:affinegraph}{gradiennya}
$0$). Bukan keduanya: $h$, karena kuadratnya.
\end{solution}

\begin{solution}{exo:g9:linfunc:2}
$f(0) = -3$; $f(4) = 5$; $f(-1) = -5$. Dan $f(x) = 9$ berarti
$2x - 3 = 9$, jadi $2x = 12$ dan $x = 6$.
\end{solution}

\begin{solution}{exo:g9:linfunc:3}
$a = \frac{15}{6} = 2.5$, jadi $f(x) = 2.5x$. Lalu $f(10) = 25$ dan
$f(-2) = -5$.
\end{solution}

\begin{solution}{exo:g9:linfunc:4}
$f$ melalui titik asal dengan \omterm{prop:g9:linfunc:affinegraph}{gradien} $2$; $g$ punya \omterm{prop:g9:linfunc:affinegraph}{gradien} yang sama
tetapi memotong sumbu-$y$ di $3$; $h$ turun dari $(0, 4)$ dengan
\omterm{prop:g9:linfunc:affinegraph}{gradien} $-1$. Grafik $f$ dan $g$ berupa \emph{garis yang \omterm{def:g6:lines:perp}{sejajar}}
(\omterm{prop:g9:linfunc:affinegraph}{gradiennya} sama, titik potongnya berbeda).
\end{solution}

\begin{solution}{exo:g9:linfunc:5}
Untuk $f$: \omterm{prop:g9:linfunc:affinegraph}{gradiennya} $\frac{11 - 5}{3 - 1} = 3$, lalu
$5 = 3 \times 1 + b$ memberi $b = 2$: jadi $f(x) = 3x + 2$.

Untuk $g$: \omterm{prop:g9:linfunc:affinegraph}{gradiennya} $\frac{0 - 4}{2 - 0} = -2$, dan $g(0) = 4$ sudah
menjadi titik potongnya: jadi $g(x) = -2x + 4$.
\end{solution}

\begin{solution}{exo:g9:linfunc:6}
Kenaikan $25\%$: $120 \times 1.25 = 150$.

Potongan $40\%$: $65 \times 0.6 = 39$.

Harga aslinya $p$ dengan $p \times 0.75 = 69$:
$p = \frac{69}{0.75} = 92$.
\end{solution}

\begin{solution}{exo:g9:linfunc:7}
\emph{1.} $f(x) = 0.05x + 10$.

\emph{2.} Paket B lebih murah ketika $25 < 0.05x + 10$, yaitu
$15 < 0.05x$, yaitu $x > 300$. Mulai dari $301$ menit ($5$ jam) tiap
bulannya, paket B menang.
\end{solution}

\begin{solution}{exo:g9:linfunc:8}
\emph{1.} Sesudah satu jam: $2000 \times 1.1 = 2200$. Sesudah dua jam:
$2200 \times 1.1 = 2420$.

\emph{2.} Kenaikan yang kedua berlaku atas $2200$, bukan atas $2000$:
tumbuh $10\%$ dua kali mengalikan dengan $1.1^2 = 1.21$, yaitu
kenaikan $21\%$, bukan $20\%$.
\end{solution}

\begin{solution}{exo:g9:linfunc:9}
\emph{1.} \omterm{prop:g9:linfunc:affinegraph}{Gradiennya}: $\frac{1 - 3}{6 - 2} = \frac{-2}{4} = -0.5$:
jadi fungsinya turun.

\emph{2.} $f(x) = -0.5x + b$ dengan $f(2) = 3$: $-1 + b = 3$, jadi
$b = 4$ dan $f(x) = -0.5x + 4$. Keluaran $0$: $-0.5x + 4 = 0$ memberi $x = 8$.
\end{solution}

\begin{solution}{exo:g9:linfunc:10}
\emph{1.} Kedua perubahannya mengalikan harganya dengan
\[
\left(1 + \frac{t}{100}\right)\left(1 - \frac{t}{100}\right)
= 1 - \frac{t^2}{10000}
\]
(kesamaan ketiga dengan $a = 1$, $b = \frac{t}{100}$).

\emph{2.} Untuk $0 < t \leq 100$, berlaku $\frac{t^2}{10000} > 0$,
jadi faktornya kurang dari $1$: harga akhirnya lebih rendah daripada
harga aslinya. Untuk $t = 10$: faktornya
$1 - \frac{100}{10000} = 0.99$, yaitu penurunan keseluruhan $1\%$.
\end{solution}

\begin{solution}{pb:g9:linfunc:1}
\textbf{1.} \omterm{prop:g9:linfunc:affinegraph}{Gradiennya}: $\frac{f(100) - f(0)}{100 - 0} =
\frac{212 - 32}{100} = 1.8$; titik potongnya $f(0) = 32$. Karena itu
\[
f(C) = 1.8\,C + 32 .
\]

\textbf{2.} $f(20) = 68\,^\circ$F (yaitu hari sejuk yang terkenal);
$f(37) = 98.6\,^\circ$F; $f(-10) = 14\,^\circ$F.

\textbf{3.} Dari $F = 1.8C + 32$:
$C = \frac{F - 32}{1.8}$. Lalu $50\,^\circ$F
$\to \frac{18}{1.8} = 10\,^\circ$C, dan $98.6\,^\circ$F
$\to \frac{66.6}{1.8} = 37\,^\circ$C.

\textbf{4.} $f(0) = 32 \neq 0$, dan $f(20) = 68$ bukan dua kali
$f(10) = 50$: jadi bukan linear. Grafiknya berupa garis lurus yang
melewatkan titik asalnya: $f$ bersifat \emph{afin}, bukan linear.

\textbf{5.} $1.8x + 32 = x$ memberi $0.8x = -32$, jadi $x = -40$:
pada empat puluh di bawah \omterm{def:g1:counting:zero}{nol}, kedua termometernya membaca $-40$. Pada
grafiknya, garis $f$ memotong garis diagonal $y = x$ tepat di situ.

\textbf{6.} $p_A(k) = 1.5k + 3$ dan $p_B(k) = 2k$. Untuk $4$ km:
$9$ euro melawan $8$: jadi B lebih murah. Untuk $8$ km: $15$ melawan
$16$: jadi A lebih murah.

\textbf{7.} $1.5k + 3 = 2k$ memberi $k = 6$: pada enam kilometer
keduanya berbiaya $12$ euro. Secara grafik, garis B (yang melalui
titik asal, lebih curam) bermula di bawah garis A lalu memotongnya di
$(6, 12)$: B menang untuk perjalanan pendek, A untuk yang panjang.

\textbf{8.} $N = 7 \times 20 - 30 = 110$ kerikan tiap menit. Dari
$82 = 7T - 30$: $T = \frac{112}{7} = 16\,^\circ$C.

\textbf{9.} $v(t) = 600 - 15t$. Bernilai $240$:
$600 - 15t = 240$, jadi $t = 24$ bulan. Bernilai $0$: $t = 40$ bulan.
Melewati $40$ bulan rumusnya menjadi negatif --- tetapi nilai sebuah
telepon berhenti di \omterm{def:g1:counting:zero}{nol}: setiap model afin punya daerah tempat model
itu masuk akal.

\textbf{10.} Dengan menyusunnya: $\times 1.25$ lalu $\times 0.8$
mengalikan dengan $1.25 \times 0.8 = 1$: jadi harganya kembali persis
ke awalnya. Bukan keajaiban: persennya bekerja atas acuan yang
berbeda (potongannya berlaku atas harga yang sudah dinaikkan), dan di
sini pengalinya kebetulan saling meniadakan --- dan kisah umumnya,
beserta
kerugian bawaannya $1 - \frac{t^2}{10\,000}$, adalah
\cref{exo:g9:linfunc:10}.

\textbf{11.} \omterm{prop:g9:linfunc:affinegraph}{Gradiennya}: $\frac{17 - 20}{30} = -0.1$ cm tiap menit;
jadi $h(t) = 20 - 0.1t$. Lilinnya mati pada $h(t) = 0$: $t = 200$
menit --- yaitu tiga jam dua puluh menit cahaya.

\textbf{12.} Yang kedua memulai $10$ km di depan (titik potongnya);
yang pertama lebih cepat ($24 > 18$ km/jam). Penyusulannya:
$24t = 10 + 18t$ memberi $6t = 10$, jadi $t = \frac53$ jam $= 1$ jam
$40$ menit, pada kilometer $24 \times \frac53 = 40$.

\textbf{13.} $f(x) + g(x) = (2 - 0.5)x + (3 + 7) = 1.5x + 10$:
jadi afin, dengan \omterm{prop:g9:linfunc:affinegraph}{gradien} berupa \emph{\omterm{def:g1:addition:def}{jumlah} kedua \omterm{prop:g9:linfunc:affinegraph}{gradiennya}} dan
titik potong berupa \omterm{def:g1:addition:def}{jumlah} kedua titik potongnya --- yang jelas
terlihat dengan mengumpulkan suku sejenis pada $(mx + p) + (m'x + p')$.

\textbf{14.} Kalau $m \neq 0$: tepat satu penyelesaian,
$x = -\frac pm$ --- garis yang tidak mendatar memotong sumbunya
sekali. Kalau $m = 0$, $p \neq 0$: tidak ada penyelesaiannya --- garis
mendatar di atas atau di bawah sumbunya tidak pernah menyentuhnya.
Kalau $m = 0$, $p = 0$: setiap $x$ adalah penyelesaiannya --- garisnya
\emph{adalah} sumbunya.

\textbf{15.} Tripel pertamanya: lajunya $\frac{9 - 5}{3 - 1} = 2$ dan
$\frac{17 - 9}{7 - 3} = 2$: sama, jadi titiknya segaris (pada
$y = 2x + 3$). Tripel keduanya: $\frac{9-5}{2} = 2$ tetapi
$\frac{16 - 9}{3} = \frac73 \neq 2$: jadi tidak segaris. \omterm{def:g9:linfunc:affine}{Fungsi afin}
persis fungsi yang laju perubahannya tetap ---
satu bilangan, yaitu \omterm{prop:g9:linfunc:affinegraph}{gradiennya}, menceritakan seluruh kisahnya; untuk
kurva, lajunya sendiri berubah dari titik ke titik, dan mengejarnya
membawa kita ke turunan, di buku Sekolah Menengah.

\end{solution}
